\section{Homogeneous Pfaffian normalization}\label{c7:pfaffian}

We now compute the scalar normalization used in the cylinder pressure.
Recall from Section~\ref{c7:finite} that, for an even list
$X=(x_1,\ldots,x_{2m})$,
\[
 \begin{split}
 d(u,v)&=u^2+\sqrt2uv+v^2,\qquad
 \mathcal P(X)=\prod_{j<k}\frac{d(x_j,x_k)}{x_j-x_k},\\
 N(X)&=\mathcal P(X)\operatorname{Pf}_{j,k}
               \frac{x_j^2-x_k^2}{d(x_j,x_k)}.
 \end{split}
\]
The coincident-site value is defined by polynomial continuation. We prove
\begin{equation}\label{j:homogeneous-normalization}
 N(1,\ldots,1)
 =\left(\frac{2p_0}{3}\right)^{m(2m-1)}
  \left(\frac23\right)^m m^{5/48+o(1)},
 \qquad p_0=2+\sqrt2.
\end{equation}
A reference Pfaffian supplies the first factor. The ratio to that
reference becomes a Gram determinant, and an even-test logarithmic
statistic supplies the remaining power.

\subsection{Confluence and the even cross Gram matrix}

Put $\eta=\pi/4$ and $x_j=e^{\eta w_j}$. The kernel of $N$ is
\[
 B_0(w)=\frac{\sinh(\eta w)}{\cosh(\eta w)+\cos(\pi/4)}.
\]
Compare it with $C(w)=\tanh(\pi w/6)$. Schur's Pfaffian identity
\cite[(1.2)]{IshikawaOkadaTagawaZeng2006} gives
\[
 \operatorname{Pf}_{j,k}\frac{y_j-y_k}{y_j+y_k}
 =\prod_{j<k}\frac{y_j-y_k}{y_j+y_k}.
\]
For clarity, multiplying the Pfaffian by $\prod_{j<k}(y_j+y_k)$
produces an alternating polynomial of Vandermonde degree, hence a
constant times $\prod_{j<k}(y_j-y_k)$. The constant is one by taking
positive $y_1\gg y_2\gg\cdots$: the limiting skew matrix with all upper
entries one has Pfaffian one by first-row induction. Taking
$y_j=e^{\pi w_j/3}$ gives the pair product for $C$. As the $w_j$ tend
to zero,
\[
 \frac{d(x_j,x_k)}{x_j-x_k}C(w_j-w_k)
 \longrightarrow
 p_0\frac{\pi/6}{\eta}=\frac{2p_0}{3}.
\]
Thus $\mathcal P(X)\operatorname{Pf}C(w_j-w_k)$ tends to
$(2p_0/3)^{m(2m-1)}$.

We next make the ratio of the two confluent Pfaffians explicit. For either
analytic odd kernel $B=B_0,C$, apply the same divided-difference
transformations to its rows and columns. At confluence the entry of
orders $r,s$ is
\[
 \frac{(-1)^s}{r!\,s!}B^{(r+s)}(0),\qquad 0\le r,s<2m.
\]
Entries with $r,s$ of the same parity vanish. After grouping the even
and odd orders, the Pfaffian is therefore, up to common signs, the
determinant of the cross matrix with $r=2a$, $s=2b+1$. With the Fourier
convention $\widehat h(k)=\int_{\mathbb R}e^{-ikw}h(w)\,dw$, this cross
entry is
\[
 -\frac{(-1)^{a+b}}{(2a)!(2b+1)!}
 \frac1{2\pi}\int_{\mathbb R}k^{2(a+b)}\widehat{B'}(k)\,dk.
\]
The signs and factorials cancel in the ratio for $B_0$ and $C$.
Consequently that ratio is the ratio of Gram determinants of
$1,k^2,\ldots,k^{2m-2}$ for the weights
$\widehat{B_0'}(k)\,dk$ and $\widehat{C'}(k)\,dk$.

Both weights are strictly positive and have all moments finite. To
compute them, put $\widetilde B_\beta(t)=\sinh t/(\cosh t+\cos\beta)$, so
\[
 B_0(w)=\widetilde B_{\pi/4}(\pi w/4),\qquad
 C(w)=\widetilde B_{\pi/2}(\pi w/6).
\]
Applying $\cos\beta+\sin\beta\,\partial_\beta$ to the transform of
$(\cosh t+\cos\beta)^{-1}$ from Lemma~\ref{j:fourier-kernel} gives
$\widehat{\widetilde B_\beta'}(\xi)
=2\pi\xi\cosh(\beta\xi)/\sinh(\pi\xi)$.
Scaling the two identities above gives
\[
 \widehat{B_0'}(k)=\frac{8k\cosh k}{\sinh4k},\qquad
 \widehat{C'}(k)=\frac{6k}{\sinh3k}.
\]
At zero these expressions use their continuous limits. For
$t=e^{-|k|}$ their ratio is
\[
 \frac{\widehat{B_0'}(k)}{\widehat{C'}(k)}
 =\frac23 e^{\phi(k)},\qquad
 \phi(k)=\log\frac{(1+t^2)(1-t^6)}{1-t^8}.
\]
Extend $\phi$ continuously by $\phi(0)=\log(3/2)$. Set $u=3k/\pi$ and
\[
 \mu(u)=\frac{2u}{\sinh\pi u},\qquad
 f(u)=\phi(\pi u/3).
\]
Use $\mu(0)=2/\pi$. The measure $\mu(u)\,du$ is a probability measure, and $f$ is real,
continuous, even, and exponentially decaying. Let $P_m$ be the
orthogonal projection in $L^2(\mu(u)\,du)$ onto
$\operatorname{span}\{1,u^2,\ldots,u^{2m-2}\}$. Changing variables in
the Gram determinants gives the exact identity
\begin{equation}\label{j:pfaffian-projection}
 N(1,\ldots,1)
 =\left(\frac{2p_0}{3}\right)^{m(2m-1)}
  \left(\frac23\right)^m
  \det(P_m e^fP_m|_{\operatorname{ran}P_m}).
\end{equation}
Here $e^f$ denotes multiplication by $e^{f(u)}$. We now evaluate this
last determinant at logarithmic order.

\subsection{A logarithmic statistic for even tests}

\begin{lemma}\label{j:even-log-statistic}
Let $\mu(u)=2u/\sinh(\pi u)$, with its continuous value at zero, and let
$P_m$ project in $L^2(\mu(u)\,du)$ onto the even polynomials of degree at
most $2m-2$. If $f:\mathbb R\to\mathbb R$ is continuous and even and
satisfies $|f(u)|\le C e^{-c|u|}$ for some $C,c>0$, then
\[
 \log\det(P_m e^fP_m|_{\operatorname{ran}P_m})
 =\frac{\log m}{2\pi}\int_{\mathbb R}f(u)\,du+o(\log m).
\]
\end{lemma}

\begin{proof}
We first describe the projection kernel. The Fourier transform of the
probability density $\mu$ is
\[
 \int_{\mathbb R}e^{isu}\mu(u)\,du=\frac1{\cosh^2(s/2)}.
\]
For example, for $s>0$ the residues at $u=in$, $n\ge1$, give
$4\sum_{n\ge1}(-1)^{n+1}ne^{-ns}$, which is the displayed function;
evenness and dominated convergence cover the remaining real $s$.
It follows that
\[
 H(z,u)=(1-iz)^{-1+iu}(1+iz)^{-1-iu}
       =\sum_{n\ge0}p_n(u)z^n
\]
generates real orthogonal polynomials of degree $n$, with leading
coefficient $2^n/n!$, parity $n$, and squared norm $n+1$. Indeed,
integrate $H(z,u)H(z',u)$ near $z=z'=0$ and use the
Fourier transform to obtain $(1-zz')^{-2}$. These are the
Meixner--Pollaczek polynomials with parameters $(1,\pi/2)$
\cite[Equation~18.23.7]{DLMF}. In particular the even members span
$\operatorname{ran}P_m$.

Uniformly for $u$ in a fixed real compact set, singularity expansion of
this generating function gives, as $j\to\infty$,
\begin{equation}\label{j:even-polynomial-asymptotic}
 p_{2j}(u)=(-1)^j\left(
 a(u)(2j)^{iu}+\overline{a(u)}(2j)^{-iu}\right)
 +O(\log j/j),
\end{equation}
where $a(-u)=\overline{a(u)}$ and
\begin{equation}\label{j:even-amplitude}
 |a(u)|^2\mu(u)=\frac1{2\pi}.
\end{equation}
Here is a direct error estimate. Subtract from $H$ its two leading
singular terms
\[
 2^{-1+iu}(1+iz)^{-1-iu}
 +2^{-1-iu}(1-iz)^{-1+iu}.
\]
Taylor expansion of the regular factors at $z=i$ and $z=-i$ bounds the
derivative of the remainder in the unit disk by
$C(1+|z-i|^{-1}+|z+i|^{-1})$, uniformly on the chosen $u$ compact.
Cauchy's formula for its derivative on the circle of radius $1-1/n$,
followed by division by $n$, bounds its coefficient by $O(\log n/n)$.
The binomial coefficients of the two
singular terms give
\[
 a(u)=2^{-1+iu}\lim_{n\to\infty}
       n^{-iu}\prod_{\ell=1}^n(1+iu/\ell).
\]
Logarithms give convergence with $O(1/n)$ error on compacts. Taking
moduli and using the sine product gives
$|a(u)|^2=\sinh(\pi u)/(4\pi u)$, proving
\eqref{j:even-amplitude} and \eqref{j:even-polynomial-asymptotic}.

The projection kernel relative to $\mu(u)\,du$ is
\[
 K_m(u,v)=\sum_{j=0}^{m-1}\frac{p_{2j}(u)p_{2j}(v)}{2j+1}.
\]
Put $D_m=\log m$. For fixed $u\ne0$,
\begin{equation}\label{j:even-kernel-peaks}
 \begin{gathered}
 \frac{K_m(u,u)}{D_m}\longrightarrow |a(u)|^2,\\
 \frac{K_m(u,\pm u+s/D_m)}{D_m}
 \longrightarrow |a(u)|^2\frac{\sin s}{s},
 \end{gathered}
\end{equation}
uniformly for $s$ in bounded intervals, with the continuous value at
$s=0$. To verify this, insert
\eqref{j:even-polynomial-asymptotic}. The summable errors contribute
$O(1)$. Comparing
$\sum_{1\le j<m}e^{ir\log(2j)}/j$ with its integral has bounded error
for $r$ in compact sets. Only $r=u-v$ near the first peak and $r=u+v$
near the second peak contribute at order $D_m$. Their two conjugate
terms yield $\sin s/s$. At $u=0$ the peaks coalesce; this single point
will not affect the integrals below.

The two peaks exhaust the leading projection norm. Indeed,
\[
 \int_{\mathbb R}|K_m(u,v)|^2\mu(v)\,dv=K_m(u,u).
\]
For fixed $u\ne0$, integrate over the two intervals
$v=\pm u+s/D_m$, $|s|\le A$. By
\eqref{j:even-kernel-peaks}, their total contribution divided by $D_m$
tends to
\[
 2|a(u)|^4\mu(u)\int_{-A}^A(\sin s/s)^2\,ds.
\]
As $A\to\infty$ this tends to $|a(u)|^2$, by
\eqref{j:even-amplitude} and
$\int_{\mathbb R}(\sin s/s)^2\,ds=\pi$. This is exactly the limit of
$K_m(u,u)/D_m$.

We also need a bound valid for all $u$:
\begin{equation}\label{j:even-kernel-bound}
 K_m(u,u)\mu(u)\le C D_m(1+|u|),\qquad m\ge2.
\end{equation}
Take $r=e^{-1/m}$. Up to an absolute factor, $K_m(u,u)$ is at most
$\sum_{n\ge0}p_n(u)^2r^n/(n+1)$. Angular averaging of the generating
function gives
\[
 \sum_{n\ge0}\frac{p_n(u)^2r^n}{n+1}
 =\frac1r\int_0^r\frac1{2\pi}
       \int_0^{2\pi}|H(\sqrt t\,e^{i\theta},u)|^2\,d\theta\,dt.
\]
The quotient $(1-iz)/(1+iz)$ has positive real part in the unit disk.
Thus
\[
 |H(z,u)|^2\le e^{\pi|u|}|1+z^2|^{-2}.
\]
The angular average of the last factor is $(1-|z|^4)^{-1}$, so the
preceding integral is at most
\[
 \frac{e^{\pi|u|}}r\int_0^r\frac{dt}{1-t^2}
 \le C e^{\pi|u|}\log m.
\]
Multiplying by $\mu(u)\le C(1+|u|)e^{-\pi|u|}$ proves
\eqref{j:even-kernel-bound}.

Dominated convergence using \eqref{j:even-kernel-peaks},
\eqref{j:even-amplitude}, and \eqref{j:even-kernel-bound} now gives
\begin{equation}\label{j:even-trace-limit}
 \frac{\operatorname{tr}(P_m fP_m)}{D_m}
 \longrightarrow\frac1{2\pi}\int_{\mathbb R}f(u)\,du.
\end{equation}
There is also a smaller fluctuation term:
\begin{equation}\label{j:even-commutator}
 \|(1-P_m)fP_m\|_2^2=o(D_m),
\end{equation}
where $\|\cdot\|_2$ is the Hilbert--Schmidt norm. The squared norm equals
\[
 \frac12\iint_{\mathbb R^2}
 (f(u)-f(v))^2|K_m(u,v)|^2\mu(u)\mu(v)\,du\,dv.
\]
For each fixed $u\ne0$, the inner integral divided by $D_m$ tends to
zero: the two peak neighborhoods exhaust the projection norm, and
continuity and evenness give $f(v)\to f(u)$ in both neighborhoods.
The diagonal bound permits dominated convergence when either variable
is in a fixed compact interval. When both variables are outside that
interval, use
$(f(u)-f(v))^2\le2f(u)^2+2f(v)^2$ and the projection identity to bound
the normalized contribution by
\[
 C\int_{|u|>A}f(u)^2(1+|u|)\,du,
\]
which tends to zero as $A\to\infty$. This proves
\eqref{j:even-commutator}. The evenness assumption is used precisely
at the second peak $v=-u$.

Finally set
\[
 L_m(s)=\log\det(P_m e^{sf}P_m|_{\operatorname{ran}P_m}),
 \qquad 0\le s\le1.
\]
This determinant is positive. If $P_{m,s}$ projects onto
$e^{sf/2}\operatorname{ran}P_m$, differentiating its Gram matrix gives
\[
 L_m'(s)=\operatorname{tr}(P_{m,s}f),\qquad
 L_m''(s)=\|(1-P_{m,s})fP_{m,s}\|_2^2.
\]
These second derivatives are bounded uniformly in $s$ by a constant
times the left side of \eqref{j:even-commutator}. To see this, regard
$U_s=e^{sf/2}P_m$ as a map from $\operatorname{ran}P_m$, and use the
isometry $U_s(U_s^*U_s)^{-1/2}$ onto
$\operatorname{ran}P_{m,s}$. The identity
\[
 (1-P_{m,s})f e^{sf/2}P_m
 =(1-P_{m,s})e^{sf/2}(1-P_m)fP_m
\]
and the bounds on $e^{sf/2}$ and $(U_s^*U_s)^{-1/2}$ give the
comparison, with a constant depending only on $\|f\|_\infty$.
Since $L_m(0)=0$ and $L_m'(0)=\operatorname{tr}(P_m fP_m)$,
\[
 L_m(1)=L_m'(0)+\int_0^1(1-s)L_m''(s)\,ds.
\]
The uniform second-derivative bound and \eqref{j:even-trace-limit}
prove the lemma.
\end{proof}

For the function in \eqref{j:pfaffian-projection},
\[
 \frac1{2\pi}\int_{\mathbb R}f(u)\,du
 =\frac3{\pi^2}\int_0^\infty\phi(k)\,dk.
\]
Expanding the three logarithms into geometric series gives
\[
 \int_0^\infty\phi(k)\,dk
 =\frac{\pi^2}{24}-\frac{\pi^2}{36}+\frac{\pi^2}{48}
 =\frac{5\pi^2}{144}.
\]
Lemma~\ref{j:even-log-statistic} therefore makes the determinant in
\eqref{j:pfaffian-projection} equal to $m^{5/48+o(1)}$. Substitution
proves \eqref{j:homogeneous-normalization}.
